\documentclass[a4paper, 10pt, conference]{ieeeconf}
\IEEEoverridecommandlockouts
\overrideIEEEmargins

\usepackage[caption=false, font=footnotesize]{subfig}
\usepackage{graphicx}
\usepackage{hyperref}
\hypersetup{breaklinks=true}
\usepackage{xurl}
\usepackage{amsmath}
\usepackage{amssymb}
\usepackage{booktabs}
\usepackage{multirow}
\usepackage{array}
\usepackage{float}

% A self-contained algorithm float. Built from base LaTeX and float
% rather than algorithm2e so the document compiles on any TeX
% installation without an extra package having to be present.
\newfloat{algorithm}{htbp}{loa}
\floatname{algorithm}{Algorithm}
\newcounter{algln}
\newcommand{\algreset}{\setcounter{algln}{0}}
\newcommand{\albox}[1]{% one numbered line, #1 already indented
  \refstepcounter{algln}%
  \noindent\makebox[1.5em][r]{\scriptsize\thealgln:}\;#1\par\vspace{0.4pt}}
\newcommand{\alplain}[1]{\noindent\makebox[1.5em][r]{}\;#1\par\vspace{0.4pt}}
\newcommand{\alind}{\hspace*{1.1em}}
\newcommand{\alkw}[1]{\textbf{#1}}
\newcommand{\alcom}[1]{{\footnotesize\itshape\;// #1}}
\newenvironment{algbody}{\par\vspace{1pt}\hrule\vspace{2pt}\algreset
  \scriptsize\setlength{\parindent}{0pt}}
  {\vspace{2pt}\hrule\vspace{1pt}}
\usepackage{stfloats}
\renewcommand{\arraystretch}{0.90}

% Let full-width floats occupy more of a page before LaTeX gives up and
% defers them, and allow a little more text on a page carrying one. The
% defaults are tuned for pages with at most one small float; this paper
% has four double-column figures and seven tables, and the stock values
% leave usable space empty at the foot of several pages.
\renewcommand{\topfraction}{0.92}
\renewcommand{\bottomfraction}{0.7}
\renewcommand{\textfraction}{0.07}
\renewcommand{\floatpagefraction}{0.82}
\renewcommand{\dbltopfraction}{0.92}
\renewcommand{\dblfloatpagefraction}{0.82}
\setcounter{topnumber}{3}
\setcounter{bottomnumber}{2}
\setcounter{totalnumber}{5}
\setcounter{dbltopnumber}{3}

% Whitespace around floats and captions. The stock separations assume a
% page with one or two floats; this paper carries fifteen, so the
% defaults spend close to a page on gaps alone. Nothing below changes
% any figure, table, number or sentence.
\setlength{\textfloatsep}{6pt plus 2pt minus 2pt}
\setlength{\dbltextfloatsep}{6pt plus 2pt minus 2pt}
\setlength{\intextsep}{5pt plus 2pt minus 2pt}
\setlength{\dblfloatsep}{7pt plus 2pt minus 3pt}
\setlength{\floatsep}{5pt plus 2pt minus 2pt}
\setlength{\abovecaptionskip}{4pt}
\setlength{\belowcaptionskip}{0pt}

% Display-equation and heading whitespace. Same principle: gaps only.
\AtBeginDocument{%
  \setlength{\abovedisplayskip}{4pt plus 2pt minus 2pt}%
  \setlength{\belowdisplayskip}{4pt plus 2pt minus 2pt}%
  \setlength{\abovedisplayshortskip}{2pt plus 1pt minus 1pt}%
  \setlength{\belowdisplayshortskip}{3pt plus 1pt minus 1pt}%
}
% Character protrusion and font expansion. This sets the same words in
% slightly less space by adjusting inter-letter and inter-word fitting,
% typically 2-4% over a paper this length. No word is removed and no
% type size changes.
% tracking is off deliberately: it letterspaces the small-caps table
% captions, which leaves numbers such as 94.2 unsearchable in the PDF.
\usepackage[protrusion=true, expansion=true, tracking=false]{microtype}

% Close the gap between reference entries. All 28 stay, in order.
\usepackage{etoolbox}
\apptocmd{\thebibliography}{%
  \setlength{\itemsep}{0pt plus 0.3pt}%
  \setlength{\parsep}{0pt}%
  \setlength{\topsep}{2pt}%
}{}{}

% Fixed so the title page does not silently change on every recompile.
\newcommand{\submissiondate}{September 2026}

\usepackage{titlesec}
\titlespacing*{\section}{0pt}{5pt plus 1pt minus 1pt}{3pt}
\titlespacing*{\subsection}{0pt}{4pt plus 1pt minus 1pt}{2pt}
\titlespacing*{\subsubsection}{0pt}{3pt plus 1pt minus 1pt}{2pt}

\title{\LARGE \bf
Autonomous Decentralised Search and Coverage with\\
Unmanned Aerial Systems Swarms
}

\author{Reddy, Mahadev Devireddy Venkata}

\begin{document}

%% ===================== COVER PAGE (excluded from page count) ===
\begin{titlepage}
    \begin{center}
        \vspace*{1cm}
        {\huge Autonomous Decentralised Search and Coverage with
        Unmanned Aerial Systems Swarms}

        \vspace{0.6cm}
        {\large By}

        \vspace{0.4cm}
        \textbf{Mahadev Reddy Devireddy Venkata}

        \vspace{1.2cm}
        {\huge\textbf{MSc Robotics Dissertation}}

        \vspace{16mm}

        \begin{center}
            
\includegraphics[scale=0.5]{logos/bristolcrest_colour.pdf}
            \hspace{8mm}
            
\includegraphics[scale=0.30]{logos/UWE_insignia.png}
        \end{center}

        \vspace{10mm}
        {\large School of Engineering Mathematics and Technology\\
        \textsc{University of Bristol}}\\
        \&\\
        {\large School of Engineering\\
        \textsc{University of the West of England}}\\

        \vspace{0.6cm}
        \begin{minipage}{10cm}
            \centering
            A MSc dissertation submitted to the University of
            Bristol and the University of the West of England in
            accordance with the requirements of the degree of
            \textsc{Master of Science in Robotics} in the Faculty
            of Engineering.
        \end{minipage}\\

        \vspace{0.5cm}
        Student ID: 2823511\\
        Supervisor: Dr Davood Asadi

        \vspace{0.6cm}
        \submissiondate
    \end{center}

    \newpage

    {\large{Declaration of own work}}

    I declare that the work in this MSc dissertation was carried
    out in accordance with the requirements of the University's
    Regulations and Code of Practice for Research Degree Programmes
    and that it has not been submitted for any other academic award.
    Except where indicated by specific reference in the text, the
    work is the candidate's own work. Work done in collaboration
    with, or with the assistance of, others, is indicated as such.
    Any views expressed in the dissertation are those of the author.

    \begin{flushright}
    Mahadev Reddy Devireddy Venkata, \submissiondate
    \end{flushright}

    \vspace{0.6cm}
    {\large{Ethics statement}}

    This project did not require ethical review as determined by my
    supervisor, Dr Davood Asadi. All work was simulation-based, with
    no human participants, animal subjects or physical hardware
    trials.

    \begin{flushright}
    Mahadev Reddy Devireddy Venkata, \submissiondate
    \end{flushright}

    \vspace{0.6cm}
    {\large{AI Usage Statement}}

    %% ------------------------------------------------------------
    %% ACTION REQUIRED BEFORE SUBMISSION.
    %% Rewrite this so it describes what you actually did, and check
    %% the wording your programme handbook requires. Do not sign the
    %% template below without reading it against your own practice.
    %% ------------------------------------------------------------
    Third-party AI tools were used in this project for debugging
    assistance with Python simulation code, for clarifying concepts
    in reinforcement learning and classical control, for the code and
    data audit reported in Section~\ref{sec:validity}, and for
    structural and language feedback on written drafts. All
    experimental results were generated by code written for this
    project.

    \vspace{0.6cm}
    {\large{Acknowledgements}}

    I thank Dr Davood Asadi for supervision throughout the project,
    and in particular for the feedback that prompted the extended
    evaluation in Section~\ref{sec:extended} and the diagnostic work
    in Section~\ref{sec:validity}.

    \vspace{0.6cm}
    {\large{Supplementary material}}

    Source code, all result files, the complete statistical output and
    the full figure set accompany this report as supplementary
    material.

\end{titlepage}

\maketitle
\thispagestyle{empty}
\pagestyle{empty}

%% ===================== ABSTRACT ================================
\begin{abstract}
UAV swarms operating without a central controller are attractive for
search and coverage in disaster response, because losing any single
agent does not end the mission. This work asks whether a
decentralised multi-agent reinforcement learning policy, executing
locally and passing no messages, can match classical reactive control
on area coverage. Four controllers are
compared on a $50\times50$ grid representing
$150\,\text{m}\times150\,\text{m}$ with six UAVs: Boids flocking,
Artificial Potential Fields (APF), a single-agent PPO hybrid, and a
parameter-shared multi-agent PPO policy in which all six UAVs learn
and none is hand-coded. Obstacle density was varied from 0\% to 30\%,
and every condition was run on 50 matched seeds across two disjoint
seed sets, giving 1{,}600 runs. Each seed reproduces an identical
obstacle layout for every controller, so comparisons use paired
Wilcoxon signed-rank tests with Holm-Bonferroni correction,
matched-pairs rank-biserial correlation and bootstrap confidence
intervals. The decentralised MARL policy exceeds the 90\%
coverage target in all eight density-by-seed-set cells by mean, and
significantly so in seven after correction, matching APF and beating
the two other controllers, which never reach it. It beats APF significantly in open terrain
(+4.7 and +5.0 percentage points at 0\% obstacles) and is
statistically indistinguishable from it at 10\% and 20\% on both seed
sets, while APF retains an advantage at the hardest density on the
development set ($-3.8$ points). All controllers transferred to
held-out layouts within 2.4 points, none significantly. Auditing the
apparatus produced three further findings. Changing only the reward
scheme, from per-agent to the shared team reward the multi-agent
literature warns against, costs 30.6 points, which exceeds the gap
between any two controllers compared here. Repeating training at four
seeds at 20\% density shows the one-million-step budget is
undertrained there: the across-seed spread is 7.8 points, closing to
4.2 at 1.85M steps where every seed overtakes APF, so at that density
the budget and not the method sets the measured parity. And
absolute-cell accounting shows that APF's apparent improvement with
obstacle density is a denominator artefact: it covers 640 fewer cells
at 30\% density than at 0\%, while the free area shrinks faster. The
same audit retracts a figure from an earlier draft of this work whose
checkpoint could not be reproduced, and reverses the conclusion drawn
from it.
\end{abstract}

%% ===================== INTRODUCTION ============================
\section{INTRODUCTION}\label{sec:intro}

\subsection{Motivation}

Covering an unknown area with a team of unmanned aerial vehicles is a
long-standing problem with immediate use in disaster response,
wildfire monitoring and post-event damage assessment~\cite{erdelj2017}.
After an earthquake or flood, searching the affected ground quickly is
what determines whether survivors are found. One UAV is too slow and
too fragile. A centrally controlled team is faster, but the controller
becomes a single point of failure, and in a collapsed or contested
environment the communication link that centralisation depends on is
exactly what fails first. Decentralised swarms avoid that by letting
each agent decide from what it senses locally, so the swarm keeps
working when agents drop out~\cite{dorigo2021}.

\subsection{Problem Statement}

A decentralised swarm must cover as much of an unknown, obstructed
area as possible within a fixed flight budget, using only local
sensing and without exchanging messages. Two families of
solution exist. Reactive controllers such as Boids and Artificial
Potential Fields are cheap, need no training, and are the de facto
baselines. Learned policies promise better coordination but face a
scaling barrier: a joint action space over $N$ agents with $|A|$
actions each grows as $|A|^N$, which is intractable beyond a handful
of agents. Whether a learned decentralised policy can actually beat a
well-tuned reactive controller on this task, under matched conditions
and with evidence strong enough to believe, is unresolved.

\subsection{Research Gap}

Three gaps in the existing work motivated this project. Most studies
of decentralised coverage test in empty or lightly obstructed
environments~\cite{zhu2023}, so it is unclear how methods behave as
the environment hardens. Almost all report coverage percentage and
stop, though battery use, collision risk and computational cost matter
just as much once a system leaves simulation. And least discussed, the
evidentiary standard is low: of the eight comparable studies in
Section~\ref{sec:related}, none reports a significance test, an effect
size, or more than ten replications per condition. A learned method
reported as beating a baseline by two points, from a single run,
against a classical controller whose parameters are not stated, is not
a result a reader can act on.

\subsection{Research Questions}

\textbf{RQ1.} Can a decentralised multi-agent PPO policy, in
which every UAV learns from local observations and sends no messages,
reach the 90\% coverage criterion across obstacle densities from 0\%
to 30\%?

\textbf{RQ2.} How does it compare against tuned Boids and APF
controllers under matched obstacle layouts, and does the comparison
survive correction for multiple testing?

\textbf{RQ3.} Do the controllers transfer to obstacle layouts unseen
during development, and how do they compare on the secondary measures
that determine deployability: energy, collision risk, computational
cost, load balance and swarm-size scaling?

\subsection{Contributions}

\begin{enumerate}
    \item A four-way comparison of Boids, APF, a single-agent PPO
    hybrid and a parameter-shared multi-agent PPO policy on identical
    matched obstacle layouts, with 50 replications per condition,
    paired significance tests, effect sizes and bootstrap intervals.
    None of the comparable studies surveyed reports the last three.
    \item A working decentralised MARL coverage system that
    matches or beats a tuned APF baseline, together with the two
    design decisions that make it learnable: a per-agent reward that
    repairs credit assignment, and a rollout length short enough to
    give PPO sufficient gradient updates.
    \item A direct measurement of the credit-assignment problem the
    multi-agent literature describes. Holding network, budget, seed
    and dynamics fixed and changing only the reward scheme costs 30.6
    percentage points of coverage, which is larger than the gap
    between any two controllers compared here.
    \item Reporting of variance across training runs, which the
    surveyed literature omits, and evidence that both the variance and
    the parity with APF at intermediate density are artefacts of the
    training budget.
    \item Evidence that percentage coverage over free cells is a
    misleading metric when free area varies with obstacle density, and
    absolute-cell accounting that corrects it.
    \item A validation audit of the apparatus itself, with three
    repaired defects, one retracted figure and one withdrawn
    experiment documented in Section~\ref{sec:validity}.
\end{enumerate}

%% ===================== RELATED WORK ============================
\section{RELATED WORK}\label{sec:related}

\textbf{Bio-inspired flocking.} Reynolds~\cite{reynolds1987}
introduced Boids as a distributed model in which each agent follows
local rules computed from neighbours inside a perception radius, with
no global state and no central controller. It is the origin of the
claim that coherent group behaviour emerges from purely local rules,
which is the claim decentralised reinforcement learning also makes. It
was built for visual plausibility, not coverage: unmodified, its
cohesion rule is harmful, since agents that stay together sweep the
same cells, so any coverage use needs a dispersion term. Recent
surveys~\cite{xu2022,alqudsi2024,arshid2025} treat bio-inspired
methods as the standard baseline for learned ones.
Shang et al.~\cite{shang2025} come closest to a bio-inspired coverage
method, though their objective is exploration efficiency, not coverage
at a fixed budget.

\textbf{Artificial potential fields.} Khatib~\cite{khatib1986} treats
the robot as a particle under attraction to its goal and repulsion
from obstacles. It is fully local, needs no map and costs constant
time per step, which makes it the natural classical comparator for a
learned local policy; Choset~\cite{choset2001} places it in the wider
coverage literature. Its known weakness is the local minimum in
concave geometry. Which escape remedy is used strongly determines
measured performance and papers routinely under-specify it, so an APF
baseline with a stated escape rule and a published sweep is a harder
comparator than the literature usually offers. Hickling et al.~\cite{hickling2025} show how the
field now treats APF, using an APF-derived reward to train a TD3
policy rather than using APF as a controller. Du~\cite{du2023} gives a
classical planning comparator using an enhanced A$^*$.

\textbf{MARL for coverage, and the credit-assignment problem.}
Schulman et al.~\cite{schulman2017} introduced PPO, now the default
policy-gradient method. Fusco et al.~\cite{fusco2025} state the
central difficulty for the multi-agent case: the joint action space
grows as $|A|^N$, which they call impractical for any significant
number of agents, and full decentralisation scales but produces
suboptimal actions and converges poorly under high partial
observability. They report training time rising from 0.065\,s at
$16\times10^3$ cells to 3{,}354\,s at $64\times10^3$.

Two standard responses exist and this project used both in turn. The
first is to sidestep the joint action space with hierarchy or
assignment, which is what most empirical work does, and each such
choice buys tractability by reintroducing something central.
Yu et al.~\cite{yu2025} use event-triggered hierarchical RL, where a
high-level module assigns regions and low-level policies execute
inside them; the assignment is a central decision, so the result
speaks to hierarchical coordination rather than to what independent
agents achieve. Leong et al.~\cite{leong2024} learn selection over a
discretised set of candidate viewpoints, which bounds the action space
but presupposes a generator with global knowledge of the map.
Xie et al.~\cite{xie2025} decompose and allocate cells centrally,
which is efficient and is not decentralised in the sense at issue
here. None of the three establishes whether agents acting on local
input alone can cover an unknown area, because in each the hardest
coordination decision is taken somewhere the agents are not. The
second response is \emph{parameter sharing}, in which one policy is
applied independently to every agent so the action space stays
$\text{Discrete}(|A|)$ while all agents learn. It is the only one of
the two that leaves the question intact, and it is the route taken
here.

Parameter sharing alone is not sufficient, because a shared policy
trained on a shared team reward suffers the credit-assignment problem:
an agent's reward is dominated by the behaviour of teammates it does
not control, so the correlation between its own action and its own
return is small and it learns little. The pathology is well documented
under the name of the lazy agent, and the standard responses are
difference rewards and counterfactual baselines, as in
COMA~\cite{foerster2018}, or value decomposition, as in
VDN~\cite{sunehag2018} and QMIX~\cite{rashid2018}. This literature is
directly relevant, because Section~\ref{sec:hybrid_ablation} shows the
earlier hybrid failed for exactly this reason, and
Section~\ref{sec:reward} adopts the simplest effective remedy: each
agent is rewarded for the cells it is first to visit.

\textbf{Evaluation practice, and why numerical comparison is not
available.} Replication counts across these papers are typically under
ten where stated at all. None reports a hypothesis test between its
method and its baselines, none reports an effect size, and train/test
separation over layouts is rare. A like-for-like coverage comparison
is therefore not possible. Fusco et al.\ report no coverage percentage
at all: their environment terminates on full exploration, so coverage
is 100\% whenever an episode completes, and their metrics are energy,
mission time and reward. The problem recurs because papers measure
different quantities, use different termination rules so fixed-budget
and time-to-completion framings are inverses, and differ in scale by
more than an order of magnitude while coverage is strongly sensitive
to the agent-to-area ratio. Table~\ref{tab:related} therefore compares
experimental design and evidentiary standard rather than headline
numbers.

\begin{table}[htb]
\centering
\caption{Experimental Design Compared Against Related Work.
n/s = not stated.}
\label{tab:related}
\scriptsize
\setlength{\tabcolsep}{3pt}
\begin{tabular}{llccc}
\toprule
\textbf{Study} & \textbf{Coverage metric} & \textbf{Dens.} &
\textbf{Seeds} & \textbf{Sig.\,/\,ES} \\
\midrule
Reynolds~\cite{reynolds1987} & none (animation) & no & --- & no \\
Fusco~\cite{fusco2025}       & none; energy, time & no & 1 & no \\
Yu~\cite{yu2025}             & 3D region coverage & no & n/s & no \\
Leong~\cite{leong2024}       & viewpoint coverage & no & n/s & no \\
Xie~\cite{xie2025}           & reduction in grids & no & n/s & no \\
Shang~\cite{shang2025}       & exploration effic. & no & n/s & no \\
Hickling~\cite{hickling2025} & distance, makespan & no & n/s & no \\
Du~\cite{du2023}             & path length & no & n/s & no \\
\midrule
\textbf{This work} & \textbf{\% free cells and absolute cells} &
\textbf{4} & \textbf{50} & \textbf{yes} \\
\bottomrule
\end{tabular}
\end{table}

%% ===================== METHODOLOGY =============================
\section{METHODOLOGY}\label{sec:methodology}

\subsection{System Architecture}\label{sec:arch}

Fig.~\ref{fig:sysarch} shows how the pieces fit together. A single
world holds the obstacle map, the shared visited map and the six UAV
positions. It advances only when an action has been supplied for every
agent, so all six move simultaneously within a timestep. The same
world class serves all four controllers, and the same seed produces
the same obstacle layout for each, which is what makes the comparison
matched rather than merely parallel.

\begin{figure*}[t]
    \centering
    \includegraphics[width=0.84\textwidth]{figures/results/system_architecture.png}
    \caption{System architecture. One world serves all four controllers
    on identical seeded layouts, advancing only once all six actions
    are supplied.}
    \label{fig:sysarch}
\end{figure*}

\subsection{Environment Definition}\label{sec:env}

A $50\times50$ discrete grid built in Python with Mesa~\cite{mesa}. At
3~m per cell the grid covers $150\,\text{m}\times150\,\text{m}$,
roughly a residential block. Six UAVs start spread around a shared
HOME point at grid centre $(25,25)$, eight cells out on the four
cardinals and $(\pm5,\pm5)$ on two diagonals, reflecting deployment
from one launch site; a five-cell safe zone around HOME is kept clear
of obstacles.

Early versions scattered single-cell obstacles at random, producing
environments nothing like a real search area. The generator was
rewritten to place rectangular building blocks three to eight cells on
a side, with walls and L- and U-shaped structures, at four densities:
0\%, 10\%, 20\% and 30\% of grid area. Fig.~\ref{fig:arena} shows one
such layout at 20\%, with the swarm part-way through an episode.

\begin{figure}[htb]
    \centering
    \includegraphics[width=\linewidth]{figures/methodology/scenario_arena.png}
    \caption{The task. Six UAVs sweep a
    $150\,\text{m}\times150\,\text{m}$ block at 20\% density, at step
    110 of 700 with 30\% covered. Layout and paths are both real: the
    deployed policy on a layout the experiments ran. Paths zigzag
    because actions are sampled, not greedy
    (Section~\ref{sec:eval}). Each agent sees only a window like the
    one drawn, and nothing joins one agent to another.}
    \label{fig:arena}
\end{figure}

Each episode runs 700 timesteps for every controller. Read against 3~m
cells at one cell per step, that scales to roughly 3~m/s and a
12-minute mission, which sits inside the endurance of a typical
multirotor. That scaling is an order-of-magnitude check and no more.
The model is purely kinematic, with no inertia, turn radius,
acceleration limit or wind, and diagonal moves cover $\sqrt{2}$ cells
per timestep, so implied speed varies with heading. The figures show
the episode length is not absurd; they do not support a flight-time
claim.
Coverage is the fraction of non-obstacle cells visited at least once,
\begin{equation}\label{eq:coverage}
    C = \frac{|\mathcal{V}|}{|\mathcal{F}|} \times 100\%
\end{equation}
with $\mathcal{V}$ the visited free cells and $\mathcal{F}$ all free
cells. Section~\ref{sec:absolute} reports $|\mathcal{V}|$ directly
alongside $C$, because the denominator changes with density.

\subsection{Justification of the 90\% Target}\label{sec:threshold}

The criterion comes from search and rescue planning, where the
governing quantity is probability of detection, a function of a
coverage factor $C_{f} = n_{s}dW/A$ for $n_s$ searchers travelling
distance $d$ with sweep width $W$ over area $A$.
Perkins~\cite{perkins2013} states that a probability of detection of
90\% requires $C_f = 1.2$, and records that standard tactics reach
63\%, 72\% and 87\%. None reaches 90\%, so the value is a recognised
and demanding planning target; the same framework underpins the US
National Search and Rescue Supplement~\cite{nss2018}. The construct
differs from the one measured here, since Eq.~\ref{eq:coverage} counts
cell entry with a sensor that never fails and is therefore an upper
bound on field performance. The criterion was fixed before any
experiment ran and is not adjusted afterwards.

\subsection{Observation Space}\label{sec:obs}

Each agent senses a $7\times7$ window centred on itself, a detection
radius of three cells or 9~m, realistic for a depth sensor such as the
Intel RealSense L515. The window encodes obstacles as $-1.0$,
unvisited free cells as $0.0$, visited cells as $1.0$ and other agents
as $2.0$, giving 49 values, plus two for normalised position
(Fig.~\ref{fig:obs_window}). Writing $p_i=(r_i,c_i)$ for agent $i$'s
cell and $\delta\in\{-3..3\}^2$ for an offset within the window,
\begin{equation}\label{eq:obs}
o_i[\delta]=
\begin{cases}
-1 & p_i{+}\delta \text{ obstacle or off-grid}\\
\;\,2 & \text{occupied by another agent}\\
\;\,1 & \text{visited by any agent}\\
\;\,0 & \text{free and unvisited}
\end{cases}
\end{equation}
with $o_i \in \mathbb{R}^{60}$ after appending
$(r_i/G, c_i/G)$ and the nine memory features of
Section~\ref{sec:memory}. Cells outside the grid are also $-1$, so an
agent cannot distinguish the world boundary from an obstacle.

The visited channel is the swarm's shared record, not the agent's
private one: a cell reads $1.0$ once any agent has occupied it. An
agent can tell from its window whether the ground ahead was already
swept, but not by whom or when. This is indirect coordination through
a shared trace, and on a real platform it would come from a shared
occupancy grid or from visible evidence of a prior pass. No agent
sends or receives a message, so the property under test is
decentralised control without message passing; complete informational
isolation is a stronger condition and this design does not meet it.
The channel is identical for Boids and APF, so it advantages no one.

\begin{figure}[htb]
    \centering
    \includegraphics[width=0.86\linewidth]{figures/methodology/fig_uav_observation_window.png}
    \caption{The $7\times7$ observation window. Values are the same
    for every controller; only the learned policy additionally
    receives the nine memory features of Section~\ref{sec:memory}.}
    \label{fig:obs_window}
\end{figure}

\subsubsection{Onboard memory features}\label{sec:memory}

The sensor window gives no gradient towards unexplored ground once
the immediate neighbourhood is swept, and a policy without one
performs no better than a random walk. Each agent therefore maintains
its own belief map of cells it has sensed, and derives nine features
from it: four quadrant scores for remaining unexplored area, a
normalised vector and distance to the nearest cell it believes
unexplored, the fraction of the arena it has sensed, and the fraction
of the episode remaining. Two of the nine, the sensed fraction and the
time remaining, are internal to the agent. The other seven are derived
from a target set that combines what this agent has sensed with the
same shared visited record the window already carries, so they extend
its range but introduce no mechanism the window does not already use.
They correspond to onboard map memory that any real UAV carries. The
observation is 60 values in total. Section~\ref{sec:ablation} measures
what these features are worth, because they are information the
classical controllers do not have, and the comparison against those
controllers is therefore not one of equal information.

\subsection{Action Space}\label{sec:action}

Each UAV moves one cell per timestep in one of eight directions, a
first-order kinematic model with no inertia, giving
$\text{Discrete}(8)$ per agent. Obstacle avoidance happens at the
planning layer: cells containing obstacles are removed from the legal
move set and each controller chooses among the remainder. Inertia,
wind and battery drain beyond the mission budget are not modelled.

\subsection{Reward Function and Gain Selection}\label{sec:reward}

Agent $i$ receives
\begin{equation}\label{eq:reward}
    r_{i,t} = w_{cov}\,\Delta v_{i,t} - w_{time}
              + \mathbb{1}[\text{done}]\;w_{team}\,C_T/100
\end{equation}
where $\Delta v_{i,t}$ counts the cells \emph{agent $i$ was first to
visit} at step $t$, and $C_T$ is final swarm coverage. Rewarding
own-first-visits rather than the global new-cell count is what makes
the signal a function of the agent's own action;
Section~\ref{sec:hybrid_ablation} shows what happens without it.

The three gains were set by scale reasoning rather than search, and
the reasoning is worth stating because the relative magnitudes matter
more than the absolute values (Fig.~\ref{fig:reward}).

\begin{figure}[htb]
    \centering
    \includegraphics[width=\linewidth]{figures/results/reward_structure.png}
    \caption{What each reward term contributes to one agent over one
    episode, on a symmetric log scale. The coverage term dominates by
    more than an order of magnitude; the other two shape behaviour
    without competing with it.}
    \label{fig:reward}
\end{figure}

$w_{cov}=3.0$ and $w_{time}=0.05$ set the exchange rate between
finding new ground and spending time: one new cell is worth 60 steps
of the time penalty. This must be large, because an agent that trades
a new cell for a shorter path defeats the objective; but it must be
finite, or a policy that stalls indefinitely near unexplored ground
incurs no cost. Sixty steps is roughly a tenth of the episode, which
makes loitering expensive without making the penalty compete with the
coverage signal. Over a full episode the coverage term accumulates to
roughly $750$ against $-35$ from the time penalty, so the ratio is
deliberately lopsided.

$w_{team}=2.0$ applies once, at termination, scaled by final swarm
coverage. Its purpose is to retain a shared interest so that agents do
not compete for the same cells, and its magnitude is set so that it
cannot dominate: at most $2.0$ against an accumulated $\approx 750$,
about $0.3\%$ of return. It is a tie-breaker between otherwise
equivalent policies, not a driver. Setting it much larger reintroduces
the credit-assignment problem the per-agent term exists to solve,
since a large shared terminal payment is again mostly determined by
teammates.

\subsection{PPO Architecture}\label{sec:ppo}

All six UAVs are driven by one shared policy acting on each agent's
own local observation (Fig.~\ref{fig:marlarch}). No agent is
hand-coded and no agent sends a message.

This is centralised training with decentralised execution. Training
pools all six agents' experience into one gradient update, which could
not be done onboard. Execution is local: each agent runs its own copy
of the network on its own $7\times7$ observation, nothing is
transmitted and no controller is in the loop. A deployed swarm needs
the shared network installed before flight and is independent after.
Parameter
sharing keeps the action space at $\text{Discrete}(8)$ rather than
$8^6 = 262{,}144$, which is the scaling wall
Fusco et al.~\cite{fusco2025} describe and which a naive joint-action
formulation hit directly during this project: the policy collapsed
onto a single repeated action and never recovered.

\begin{figure}[htb]
    \centering
    \includegraphics[width=0.90\linewidth]{figures/results/marl_architecture.png}
    \caption{Parameter-shared multi-agent PPO. One network produces an
    action for each agent from that agent's own observation, and each
    agent's experience is a separate sample for the same update.}
    \label{fig:marlarch}
\end{figure}

The network is a $[256,256]$ MLP with separate policy and value heads,
the Stable-Baselines3~\cite{sb3} \texttt{MlpPolicy} default. Six
agents in each of eight parallel worlds present 48 agent-slots to the
learner, so one policy update draws on all six agents' experience
simultaneously.

\subsection{Training Procedure}\label{sec:training}

Algorithm~\ref{alg:train} states the loop. The centralised step is
line~\ref{ln:update}, which pools all $NW$ agent-slots into one
gradient update; nothing else crosses between agents, and at execution
only lines~\ref{ln:obs}--\ref{ln:act} run, on each agent's own copy
of $\pi_\theta$.

\begin{algorithm}[htb]
\caption{Parameter-shared MARL training (CTDE)}
\label{alg:train}
\begin{algbody}
\alplain{\alkw{Input:} density $d$, budget $T$, seed pool
  $\mathcal{S}_A$, $W$ worlds, $N$ agents}
\alplain{\alkw{Output:} one shared policy $\pi_\theta$}
\vspace{1pt}\hrule\vspace{2pt}
\albox{initialise $\pi_\theta$, a $[256,256]$ MLP shared by all agents}
\albox{$w_1..w_W \leftarrow$ worlds at density $d$, seeds from
  $\mathcal{S}_A$}
\albox{$t \leftarrow 0$}
\albox{\alkw{while} $t < T$ \alkw{do}}
\albox{\alind $\mathcal{B} \leftarrow \emptyset$}
\albox{\alind \alkw{for} $n_{steps}$ \alkw{do}}
\albox{\alind\alind \alkw{foreach} world $w$, agent $i \in 1..N$
  \alkw{do}}
\albox{\label{ln:obs}\alind\alind\alind $o_i \leftarrow
  \textsc{Observe}(w,i)$ \alcom{Eq.\,\ref{eq:obs}}}
\albox{\label{ln:act}\alind\alind\alind $a_i \sim
  \pi_\theta(\cdot \mid o_i)$ \alcom{sampled, not greedy}}
\albox{\alind\alind $\mathbf{r},\,done \leftarrow
  \textsc{Step}(w,\mathbf{a})$ \alcom{all $N$ move together}}
\albox{\alind\alind $\mathcal{B} \leftarrow \mathcal{B} \cup
  \{(o_i,a_i,r_i)\}_{i=1}^{N}$ \alcom{$r_i$ own-first, Eq.\,\ref{eq:reward}}}
\albox{\alind\alind \alkw{if} $done$ \alkw{then} reset $w$ with a fresh
  seed}
\albox{\label{ln:update}\alind $\theta \leftarrow
  \textsc{PPO-Update}(\theta, \mathcal{B})$
  \alcom{pools all $NW$ slots}}
\albox{\alind $t \leftarrow t + N W n_{steps}$}
\albox{\alkw{return} $\pi_\theta$}
\end{algbody}
\end{algorithm}

Hyperparameters:
$n_{steps}=64$ per slot (3{,}072 samples per update), batch 512, four
epochs, $\gamma=0.99$, $\lambda=0.95$, clip 0.2, entropy 0.005,
learning rate $5\times10^{-4}$, one million agent-steps per density,
$W=8$ and $N=6$. Training layouts are drawn from Seed Set A only.


Rollout length was the decisive hyperparameter. With 48 agent-slots,
$n_{steps}=256$ yields 12{,}288 samples per update but only about
forty gradient updates per 500k steps, far too few for PPO to
converge: coverage stalls near the random-walk level of 50\% and does
not improve with further training. Shortening the rollout to 64 gives
roughly 160 updates over the same sample budget, and is what makes the
task learnable. The effective number of updates falls as the slot
count rises, which anyone scaling parameter sharing to many agents
has to compensate for.

\textbf{Training cost.} On four CPU threads of an Apple M1 the
implementation sustains about 11{,}400 agent-steps per second, so one
million takes 1.5 minutes and the four-density set about six, with no
GPU. Fusco et al.~\cite{fusco2025} report 3{,}354\,s for one run at
$64\times10^3$ cells; the environments differ too much for a direct
comparison, but parameter sharing evidently keeps a decentralised
policy within minutes of laptop compute at this scale. The classical
controllers need no training at all, and that asymmetry belongs in any
deployment decision.

\subsection{Evaluation Procedure}\label{sec:eval}

Sampling at line~\ref{ln:act} is a property of the deployed
controller, not exploration noise left switched on. Six agents share
one network, so a deterministic argmax makes every agent in an
identical local situation choose the identical action, collapsing the
swarm into one repeated move and costing almost ninety points of
coverage when measured directly. Sampling breaks the symmetry between
otherwise indistinguishable agents, and is the mechanism by which one
shared network produces six behaviours. A fielded swarm would run the
same stochastic policy. Sampling is seeded from the episode seed, so a
seed reproduces the identical episode.

\subsection{Classical-Controller Implementation}\label{sec:classical}

\textbf{Boids}~\cite{reynolds1987} scores neighbouring cells as
$s = w_{cov}c_{cov} + w_{sep}c_{sep} + \mathcal{U}(-0.1,0.1)$, with
$c_{cov}$ equal to 2 for an unvisited cell and $-1$ for a visited one,
and $c_{sep}$ the normalised distance to the nearest agent. Agents
move sequentially within a timestep, so each sees where the others
have gone, which produces spatial separation without communication.

\textbf{APF}~\cite{khatib1986} scores a candidate cell as
\begin{equation}\label{eq:apf_score}
    s = c_{cov}
      - w_{rep} \!\!\sum_{o \in \mathcal{O}_{n}} \frac{k_{rep}}{d_o^2}
      + c_{sep} + \mathcal{U}(-0.1, 0.1)
\end{equation}
over obstacles within radius $r_{rep}$, with $c_{cov}$ of 5.0 for a
cell unvisited by any agent, 1.0 for one absent from the agent's own
memory and $-2.0$ otherwise, and $w_{rep}=0.3$. The repulsion term of
Eq.~\ref{eq:apf_score} is what Section~\ref{sec:sweep} sweeps. An
escape rule fires a random move to an unvisited neighbour after three
stalled steps,
which is what prevents the local-minimum collapse the method is known
for. Rather than assert that $k_{rep}=1.5$ and $r_{rep}=2$ are well
chosen, Section~\ref{sec:sweep} reports the coverage surface over
both.

\subsection{Experimental Design and Statistical Methodology}
\label{sec:stats}

The experiment crosses four controllers with four densities and two
seed sets: 32 conditions, 50 seeds each, 1{,}600 runs. Seed Set~A
(0--49) was used for development and training, Seed Set~B (50--99)
reserved and used only at the end. Every controller sees the identical
50 layouts within a cell, as Algorithm~\ref{alg:eval} states: each
passes its seed to the same obstacle generator, and the layouts are
regenerated and compared cell by cell before any test runs, which is
what licenses the paired design. The implementation is Python~3.13
with Mesa~\cite{mesa}, Stable-Baselines3~\cite{sb3} and Gymnasium; the
bootstrap is seeded at 20260729, and code, results and checkpoints
accompany this report.

\begin{algorithm}[htb]
\caption{Matched-layout evaluation}
\label{alg:eval}
\begin{algbody}
\alplain{\alkw{Input:} controllers $\mathcal{C}$, densities
  $\mathcal{D}$, seed sets $\mathcal{S}_A,\mathcal{S}_B$}
\vspace{1pt}\hrule\vspace{2pt}
\albox{\alkw{foreach} $d \in \mathcal{D}$, $\mathcal{S} \in
  \{\mathcal{S}_A,\mathcal{S}_B\}$, seed $s \in \mathcal{S}$ \alkw{do}}
\albox{\alind $M_s \leftarrow \textsc{GenObstacles}(d,s)$;
  \alkw{assert} identical for every $c \in \mathcal{C}$}
\albox{\alind $C[c,s] \leftarrow \textsc{Run}(c,M_s,700)$
  \alkw{for each} $c \in \mathcal{C}$}
\albox{\alkw{foreach} pair $(c_1,c_2)$, each $(d,\mathcal{S})$
  \alkw{do}}
\albox{\alind Wilcoxon on $\{C[c_1,s]-C[c_2,s]\}_s$; $r_{rb}$,
  Cliff's $\delta$, bootstrap CI}
\albox{Holm--Bonferroni within each comparison family}
\end{algbody}
\end{algorithm}

Three measures in Table~\ref{tab:extended} each admit more than one
reasonable definition, and every controller is measured on the same
one. With $u_i$ the distinct cells agent $i$ occupied,
\begin{equation}\label{eq:secondary}
L=\ell\!\!\sum_{i=1}^{N}\sum_{t}\lVert p_{i,t}-p_{i,t-1}\rVert_2,
\quad
\mathcal{G}=\frac{2\sum_{k} k\,u_{(k)}}{N\sum_{k}u_{k}}-\frac{N{+}1}{N}
\end{equation}
with $\ell=3$~m per cell, $u_{(k)}$ the $k$th smallest $u_i$, and
$\mathcal{G}=0$ perfectly even. A near-miss is a pair within one cell
by Euclidean distance, counted before each move.

Twenty-seven of the 32 condition-level coverage distributions depart
from normality under Shapiro-Wilk at $\alpha = 0.05$, so rank-based
tests are primary. The design is matched, and this was verified rather
than assumed: for each seed the obstacle map was regenerated
independently through the Mesa generator and through both learned
environments and compared cell by cell. The maps are identical for
every sampled seed at every density, so observation $i$ in each sample
refers to the same environment and the paired Wilcoxon signed-rank
test is correct. Mann-Whitney $U$~\cite{mann1947} and Welch's $t$ are
reported as robustness checks. Holm-Bonferroni
correction~\cite{holm1979} is applied within each family. Effect sizes
are matched-pairs rank-biserial correlation $r_{rb}$ and Cliff's
$\delta$~\cite{cliff1993}; intervals are bootstrap percentile
intervals from 10{,}000 resamples~\cite{efron1979}, and all error bars
in this report are 95\% confidence intervals on the mean rather than
standard deviations.

%% ===================== RESULTS =================================
\section{RESULTS}\label{sec:results}

Table~\ref{tab:main_results} reports final coverage for every
controller at every obstacle density on both seed sets, with the
number of runs reaching the 90\% target in parentheses.

\begin{table}[htb]
\centering
\caption{Final Coverage (\%), Mean $\pm$ SD over 50 Seeds, with Runs
Reaching the 90\% Target in Parentheses}
\label{tab:main_results}
\scriptsize
\setlength{\tabcolsep}{0.8pt}
\resizebox{\columnwidth}{!}{%
\begin{tabular}{@{}l@{\hspace{2pt}}lcccc@{}}
\toprule
\textbf{Set} & \textbf{Controller} & \textbf{0\%} & \textbf{10\%} &
\textbf{20\%} & \textbf{30\%} \\
\midrule
\multirow{4}{*}{A}
 & Boids  & $87.9\pm4.5$\,(17) & $86.7\pm4.8$\,(12) & $85.7\pm5.6$\,(12) & $84.6\pm6.6$\,(9) \\
 & APF    & $91.4\pm3.6$\,(36) & $92.5\pm3.9$\,(39) & $94.2\pm3.5$\,(42) & $\mathbf{94.7\pm2.7}$\,(47) \\
 & Hybrid & $85.8\pm5.7$\,(13) & $85.0\pm5.5$\,(9)  & $84.8\pm7.0$\,(13) & $83.8\pm6.1$\,(10) \\
 & MARL   & $\mathbf{96.1\pm1.5}$\,(50) & $\mathbf{93.7\pm2.5}$\,(45) & $93.5\pm3.7$\,(41) & $90.8\pm5.7$\,(35) \\
\midrule
\multirow{4}{*}{B}
 & Boids      & $86.6\pm4.6$\,(12) & $86.8\pm6.4$\,(18) & $87.5\pm4.9$\,(17) & $83.6\pm6.7$\,(7) \\
 & APF        & $90.5\pm4.1$\,(27) & $93.0\pm3.8$\,(41) & $\mathbf{94.5\pm2.8}$\,(47) & $\mathbf{94.0\pm4.4}$\,(42) \\
 & Hybrid     & $85.6\pm5.2$\,(11) & $85.9\pm4.7$\,(8)  & $84.4\pm6.3$\,(9)  & $82.2\pm6.8$\,(5) \\
 & MARL   & $\mathbf{95.5\pm1.8}$\,(49) & $\mathbf{93.8\pm2.7}$\,(46) & $92.8\pm3.8$\,(36) & $93.1\pm3.8$\,(41) \\
\bottomrule
\end{tabular}}
\end{table}

\begin{figure*}[t]
    \centering
    \includegraphics[width=0.84\textwidth]{figures/results/four_way_coverage_ci.png}
    \caption{Coverage against obstacle density for all four
    controllers, on the development and held-out seed sets. Error bars
    are 95\% bootstrap confidence intervals on the mean. Only MARL and
    APF clear the target, crossing between 10\% and 20\% on both sets.}
    \label{fig:fourway}
\end{figure*}

\subsection{Meeting the 90\% Criterion}

Both the MARL policy and APF exceed 90\% by mean in all eight of
their cells. One-sided Wilcoxon tests against 90\%, Holm-corrected
over all 32 cells, confirm the excess is significant in seven cells
each. The exceptions are the cells where the mean sits closest to the
line: MARL at 30\% density on the development set (90.84\%, Holm
$p = 0.499$) and APF at 0\% density on the held-out set (90.46\%,
Holm $p = 1.000$). Boids falls significantly short in seven cells and
is indistinguishable from 90\% in one; the hybrid falls short in all
eight. Only the two leading controllers meet the criterion at all,
and neither dominates the other on it.

\subsection{Pairwise Comparison}\label{sec:pairwise}

\begin{table}[htb]
\centering
\caption{MARL Against Each Baseline, Paired Wilcoxon Signed-Rank,
Holm-Corrected. Difference in percentage points.}
\label{tab:stats}
\scriptsize
\setlength{\tabcolsep}{1.6pt}
\begin{tabular}{lllrcrr}
\toprule
\textbf{Set} & \textbf{Dens.} & \textbf{Comparison} & \textbf{Diff.} &
\textbf{95\% CI} & \textbf{Holm $p$} & \textbf{$r_{rb}$} \\
\midrule
\multirow{6}{*}{A}
& 0\%  & MARL vs APF    & $+4.71$  & [3.67, 5.81]        & $<$0.0001 & $+0.94$ \\
& 10\% & MARL vs APF    & $+1.21$  & [$-0.09$, 2.53]     & 0.5523    & $+0.24$ \\
& 20\% & MARL vs APF    & $-0.65$  & [$-2.02$, 0.72]     & 0.6892    & $-0.20$ \\
& 30\% & MARL vs APF    & $-3.82$  & [$-5.59$, $-2.15$]  & 0.0004    & $-0.64$ \\
& 0\%  & MARL vs Boids  & $+8.19$  & [7.00, 9.40]        & $<$0.0001 & $+1.00$ \\
& 0\%  & MARL vs Hybrid & $+10.28$ & [8.72, 11.95]       & $<$0.0001 & $+1.00$ \\
\midrule
\multirow{6}{*}{B}
& 0\%  & MARL vs APF    & $+5.02$  & [3.72, 6.31]        & $<$0.0001 & $+0.87$ \\
& 10\% & MARL vs APF    & $+0.82$  & [$-0.59$, 2.30]     & 0.9649    & $+0.11$ \\
& 20\% & MARL vs APF    & $-1.73$  & [$-2.99$, $-0.51$]  & 0.2702    & $-0.33$ \\
& 30\% & MARL vs APF    & $-0.90$  & [$-2.45$, 0.68]     & 0.7116    & $-0.24$ \\
& 20\% & MARL vs Boids  & $+5.33$  & [3.69, 7.06]        & $<$0.0001 & $+0.77$ \\
& 20\% & MARL vs Hybrid & $+8.41$  & [6.41, 10.49]       & $<$0.0001 & $+0.90$ \\
\bottomrule
\end{tabular}

\vspace{0.3em}
{\footnotesize The full 48-comparison table, including every
MARL-vs-Boids and MARL-vs-Hybrid cell and the Mann-Whitney and Welch
robustness checks, is in the supplementary output.}
\end{table}

Three findings follow. MARL beats APF in open terrain by 4.71 and
5.02 points, both highly significant with near-maximal effect sizes.
The two cannot be separated at intermediate density: at 10\% and 20\%
no comparison survives Holm correction on either seed set. APF keeps
an edge at the hardest density by 3.82 points with a medium effect,
but only on the development set; the same comparison on the held-out
set is 0.90 points and not significant, so the advantage is not robust
across layout samples. Against the two weaker baselines MARL wins
everywhere by 5.3 to 10.9 points, $r_{rb}$ 0.75 to 1.00, all sixteen
surviving correction.

\subsection{Generalisation}\label{sec:generalisation}

Seed Set~B consists of layouts distinct from Set~A, verified by
regenerating and comparing maps. Changes from A to B for the MARL
policy are $+0.60$, $-0.11$, $+0.76$ and $-2.31$ percentage points at
0--30\% density; no controller changes by more than 2.31 points and no
difference survives Holm correction. The policy was trained only on
Set~A layouts, so this is a genuine held-out test and shows no
overfitting to the training maps; at 30\% density it performs
\emph{better} on the unseen set. At 0\% density there are no obstacles,
so the two sets are the same environment and any difference there
reflects only the stochastic policy.

\subsection{Absolute Cells Against Percentage}\label{sec:absolute}


Absolute-cell accounting resolves an ambiguity the earlier iteration
of this project could not settle. APF's percentage coverage rises by 3.3 points from 0\% to
30\% obstacles, which invites the reading that structured obstacles
help it search. The absolute counts, recovered over the same 50 seeds
by multiplying each run's coverage by the free area of its layout,
contradict that. APF covers 640 fewer cells at 30\% than at 0\%, a
fall of 28\%; Boids 727 fewer, 33\%; and MARL 824 fewer, 34\%. Every
controller covers less ground as density rises, while the free area
itself shrinks from 2{,}500 to 1{,}737 cells, a fall of 31\%. APF's
percentage rises only because its absolute coverage falls more slowly
than the denominator. No controller becomes
a better searcher in cluttered terrain; APF degrades most gracefully.
Percentage coverage alone is a misleading metric across varying free
area, and reporting both is the correct practice.

Fig.~\ref{fig:traj_30} shows what that difference looks like on the
ground: three controllers sweeping one identical 20\% layout, where
APF's corridor-following and MARL's outward partition are visible as
search patterns rather than as summary statistics.

\begin{figure*}[t]
    \centering
    \includegraphics[width=0.84\textwidth]{figures/results/trajectories_four_way.png}
    \caption{Swarm trajectories on one identical layout at 20\%
    obstacle density. Boids doubles back in open areas; APF follows
    the corridors between buildings; the MARL policy sweeps
    outward and partitions the arena without communicating.
    $\star$ marks HOME. Equivalents at every density are in the
    supplementary material.}
    \label{fig:traj_30}
\end{figure*}

%% ===================== EXTENDED ANALYSIS =======================
\section{EXTENDED ANALYSIS}\label{sec:extended}

Coverage alone is a poor guide to field viability. Following
supervisory feedback, further analyses were run at the same 700-step
budget on 50 matched seeds, except the swarm-size and hybrid sweeps,
which need a run per configuration and use 15.

\begin{table}[htb]
\centering
\caption{Secondary Measures, Mean over 50 Seeds, defined at
Eq.~\ref{eq:secondary}. Near-misses are a proximity proxy, not
observed collisions. $^\dagger$Per-step cost re-measured for all four
on one machine; only the ordering is meaningful.}
\label{tab:extended}
\scriptsize
\setlength{\tabcolsep}{3pt}
\begin{tabular}{llcccc}
\toprule
\textbf{Measure} & \textbf{Controller} & \textbf{0\%} & \textbf{10\%} &
\textbf{20\%} & \textbf{30\%} \\
\midrule
\multirow{4}{*}{Path length}
 & Boids  & 14913 & 14860 & 14819 & 14747 \\
 & APF    & 14946 & 14837 & 14723 & 14582 \\
 & Hybrid & 13927 & 12744 & 12690 & 13505 \\
 & MARL   & 15240 & 14399 & 13696 & 11950 \\
\midrule
\multirow{4}{*}{Gini}
 & Boids  & 0.067 & 0.077 & 0.087 & 0.103 \\
 & APF    & 0.029 & 0.039 & 0.050 & 0.083 \\
 & Hybrid & 0.197 & 0.228 & 0.239 & 0.244 \\
 & MARL   & 0.058 & 0.051 & 0.064 & 0.103 \\
\midrule
\multirow{4}{*}{Near-misses}
 & Boids  & 1244.4 & 1088.4 & 888.9 & 701.5 \\
 & APF    & 1245.4 & 1109.4 & 938.7 & 606.6 \\
 & Hybrid & 0.8    & 0.9    & 0.7   & 0.8   \\
 & MARL   & 47.8   & 55.6   & 50.1  & 69.6  \\
\midrule
\multirow{4}{*}{ms\,/\,step$^\dagger$}
 & Boids  & 0.085 & 0.212 & 0.327 & 0.431 \\
 & APF    & 0.358 & 0.581 & 0.652 & 0.776 \\
 & Hybrid & 0.412 & 0.418 & 0.408 & 0.392 \\
 & MARL   & 0.728 & 0.889 & 0.702 & 0.683 \\
\bottomrule
\end{tabular}
\end{table}

\textbf{Load balance.} APF spreads work most evenly, MARL second,
Boids third. The hybrid is far worse, and the per-agent breakdown
explains why: at 20\% obstacles its agents contribute 0.4\%, 19.6\%,
19.7\%, 20.9\%, 19.6\% and 19.8\% of unique coverage. Agent~0, the
only one under policy control, accounts for 0.4\% against a balanced
16.7\%. MARL shows no such pathology and reaches near-APF balance with
no agent hand-coded to do so, because each is rewarded for its own
contribution.

\textbf{Proximity events.} A near-miss proxies collision risk and
counts no collisions: the grid gives a vehicle no body, so the measure
captures how often a real airframe would have had to give way. MARL
records roughly twenty times fewer than either classical controller
while covering more ground. The hybrid's near-zero figure comes from
elsewhere: read with the load-balance result, it describes its single
learned agent staying out of the way.

\textbf{Computational cost.} Boids and APF are $O(N^2)$ per timestep,
each agent checking every other; a shared policy is $O(N)$, one
forward pass each. The constants reflect that: the classical pair
rises with density, the learned ones stay flat, and at six agents MARL
is still dearest of the four, a forward pass costing more than
scanning five neighbours. The scaling favours it; at this size it buys
coverage with compute.

\textbf{Swarm size.} All four controllers were run at 3, 6, 9 and 12
agents at 20\% density over 15 seeds (Table~\ref{tab:swarm}). MARL and
the hybrid were trained at six agents, so other sizes are a transfer
test; the shared-policy architecture makes this possible at all, since
one network applies to any swarm size without retraining.

\begin{table}[htb]
\centering
\caption{Coverage and Near-Misses by Swarm Size, 20\% Obstacles,
15 Seeds. MARL and the hybrid are six-agent policies tested at other
sizes. All four columns count a near-miss identically: Euclidean
radius one cell, sampled before each step.}
\label{tab:swarm}
\scriptsize
\setlength{\tabcolsep}{3pt}
\begin{tabular}{lcccc|cccc}
\toprule
& \multicolumn{4}{c|}{Coverage (\%)} & \multicolumn{4}{c}{Near-misses} \\
$N$ & Boids & APF & Hyb. & MARL & Boids & APF & Hyb. & MARL \\
\midrule
3  & 63.5 & 68.9 & 49.4 & \textbf{73.6} & 185  & 175  & 2.2  & 5    \\
6  & 85.2 & 93.5 & 82.6 & \textbf{93.6} & 827  & 910  & 20.9 & 52   \\
9  & 93.9 & \textbf{99.0} & 92.8 & 98.9 & 2312 & 2059 & 55.7 & 200  \\
12 & 97.7 & \textbf{99.7} & 97.3 & 99.2 & 3845 & 3819 & 92.9 & 300  \\
\bottomrule
\end{tabular}
\end{table}

Two things stand out. The MARL policy leads at three agents by 4.7
points over APF and by 10.1 over Boids, which is the condition where
coordination matters most because there is least redundancy to hide a
poor sweep. It converges with APF at six and nine and falls slightly
behind at twelve, where the arena is saturated and any reasonable
controller approaches the ceiling. Returns fall off past nine agents
for every controller while near-misses rise four- to sixfold from six
to twelve (Fig.~\ref{fig:sensitivity}), so six agents sits on the
steep part of the curve where coordination differences are amplified.
Across all sizes the learned controllers record twelve to thirty-four
times fewer near-misses than the classical ones.

\begin{figure*}[t]
    \centering
    \includegraphics[width=0.84\textwidth]{figures/extended/uav_count_sensitivity_all3.png}
    \caption{Coverage and near-misses against swarm size at 20\%
    density, 15 seeds, the same runs as Table~\ref{tab:swarm}. The
    learned controllers are dashed: both were trained at six agents
    and tested outside that configuration at every other size. Log
    axis on the right, where they stay about an order of magnitude
    below the classical pair throughout.}
    \label{fig:sensitivity}
\end{figure*}

\textbf{Moving obstacles.} With obstacles relocating every 50 steps at
20\% density, APF falls from 94.0\% to 85.9\% as the moving share goes
from none to half, while Boids drops only from 85.9\% to 82.8\%. APF
needs a field consistent long enough to guide an agent through; Boids
recomputes each step. The learned policies were not retrained under
this perturbation, so this covers the classical pair only.

\subsection{Information Asymmetry: the Memory Ablation}
\label{sec:ablation}

The learned controller carries nine onboard-memory features
(Section~\ref{sec:memory}) that Boids and APF do not have. Rather than
declare this as a caveat, it was measured. An otherwise identical
policy was trained with those features zeroed, leaving it exactly the
information the classical controllers receive: a $7\times7$ sensed
window and its own position.

\begin{table}[htb]
\centering
\caption{Memory-Feature Ablation, 20\% Density, 50 Matched Seeds. The
Ablated Policy Keeps the Same 60-Dimensional Vector With the Nine
Memory Features Held at Zero}
\label{tab:ablation}
\scriptsize
\setlength{\tabcolsep}{3pt}
\begin{tabular}{lccc}
\toprule
\textbf{Configuration} & \textbf{Steps} & \textbf{Coverage (\%)} &
\textbf{$\geq$90\%} \\
\midrule
MARL, memory     & 1.00M & $93.6\pm3.7$ & 41/50 \\
MARL, memory     & 1.85M & $98.8\pm1.0$ & 50/50 \\
MARL, no memory  & 1.00M & $89.6\pm4.6$ & 31/50 \\
MARL, no memory  & 1.85M & $90.0\pm4.2$ & 29/50 \\
APF (same information) & --- & $94.2\pm3.5$ & 42/50 \\
\bottomrule
\end{tabular}
\end{table}

At a matched budget of one million agent-steps the memory features are
worth $+4.00$ percentage points (paired Wilcoxon, $p<0.0001$,
$r_{rb}=0.72$), and without them the policy reaches 89.6\% against
APF's 94.2\%. On equal information and equal training the classical
controller wins, and it still wins at 1.85M ($-4.22$ pp, $p<0.0001$).

Running both arms on to 1.85M separates two explanations that the
single-budget comparison cannot. The memory-free policy gains nothing
from the extra 850k steps ($+0.43$ pp, $p=0.83$): it has converged. The
memory policy gains $+5.26$ pp over the same interval and reaches
98.8\%, so the gap widens from 4.00 to 8.83 points
($r_{rb}=1.00$). Extra training does not close it, which rules out the
head-start explanation. What the features change is what the
architecture can represent from local sensing alone: a $7\times7$
window and a position alone do not reach the target reliably.

\subsection{Credit Assignment Measured Directly}
\label{sec:credit}

Section~\ref{sec:reward} rewards each agent for the cells it is first
to visit. The alternative that the literature warns against is the
naive shared reward, in which every agent receives the global new-cell
count. Difference rewards~\cite{foerster2018} and value
decomposition~\cite{sunehag2018},~\cite{rashid2018} exist to repair
that arrangement, and it can be run here as a controlled experiment,
since the network, budget, seed, observation and dynamics all stay
fixed and only the reward scheme moves.

The shared-reward policy reaches $63.0\pm6.3\%$ against $93.6\pm3.7\%$
for the per-agent reward, a gap of $30.60$ percentage points
($p<10^{-14}$, $r_{rb}=1.00$), and not one of its fifty runs reaches
the target. On this task credit assignment decides whether the
controller works at all: on the shared reward it loses to every
baseline. It also accounts for the hybrid in
Section~\ref{sec:hybrid_ablation} without appeal to its architecture,
since the hybrid used exactly this reward.

\subsection{Training-Seed Variance and Budget}\label{sec:seedvar}

Every MARL figure reported above comes from one training run at PPO
seed 0. Fifty evaluation seeds control layout variance; they say
nothing about variance across training runs. Three further policies
were trained at seeds 1, 2 and 3 under an identical budget and
evaluated on the same fifty layouts (Table~\ref{tab:seedvar}).

\begin{table}[htb]
\centering
\caption{Coverage by PPO Training Seed, 20\% Density, 50 Matched
Evaluation Seeds Each. APF Scores 94.2\% on the Same Layouts.
SD$_{\text{seed}}$ is Across the Four Means, not Across Runs}
\label{tab:seedvar}
\scriptsize
\setlength{\tabcolsep}{2pt}
\begin{tabular}{lccccccc}
\toprule
\textbf{Budget} & \textbf{s0} & \textbf{s1} & \textbf{s2} &
\textbf{s3} & \textbf{Mean} & \textbf{SD$_{\text{seed}}$} &
\textbf{$\geq$90\%} \\
\midrule
1.00M & 93.55 & 92.98 & 91.08 & 85.75 & 90.84 & 3.55 & 122/200 \\
1.85M & 98.81 & 98.87 & 96.81 & 94.69 & 97.30 & 1.98 & 193/200 \\
\bottomrule
\end{tabular}
\end{table}

At 1.00M the spread is wide: 7.8 points between the best and worst
seed, and seed 0 is the best of the four. The headline MARL numbers
therefore sit at the top of their distribution, not in the middle of
it. A single fresh training run at 20\% density averages 90.8\%, with
a standard deviation of 5.0 points across the 200 individual runs, and
that puts it behind APF.

The 1.85M row identifies the cause. The across-seed standard deviation
falls from 3.55 to 1.98 points, the worst seed rises to 94.7\% and
overtakes APF, and every seed clears the target on at least 44 of 50
layouts. Both the shortfall and the instability at 1.00M come from
stopping on the steep part of the learning curve. At 1.85M the same
policy beats APF by 4.61 points at 20\% density ($p<10^{-14}$,
$r_{rb}=1.00$), so what Table~\ref{tab:stats} records at that density
is the training budget talking and not a ceiling on the method.
Holding the budget fixed across every controller and density is the
right control for the main comparison, but at 20\% it understates the
learned controller.

How far that generalises is untested. The extended-budget evidence
comes from one density, and 10\% is a different regime with a smaller
starting margin and sparser structure for APF to exploit. The same
direction is plausible there, but plausible is an inference, not a
measurement. Repeating the protocol at every density is future work.

\subsection{APF Parameter Sensitivity}\label{sec:sweep}

The repulsion parameters were swept over $k_{rep} \in \{0.5, 1.0,
1.5, 2.5, 4.0\}$ and $r_{rep} \in \{1,2,3,4\}$ at every density. The
reported configuration $(1.5, 2)$ sits on a broad plateau: coverage
varies by under two points across the centre of the grid and degrades
sharply only at $r_{rep}=4$, where the repulsion radius approaches the
corridor width. APF's showing is a property of the method here, not of
a fortunate parameter choice, which is what makes it a fair
comparator. The full surface is in the supplementary output.

%% ===================== VALIDATION ==============================
\section{VALIDATION AND THREATS TO VALIDITY}\label{sec:validity}

Three defects found by auditing the experiment code were repaired
before any result above was generated, and one earlier experiment is
withdrawn. A fourth item, the ablation in
Section~\ref{sec:hybrid_ablation}, is reported as a finding rather
than a defect.

\subsection{The Hybrid Could Not Answer the Question}
\label{sec:hybrid_ablation}

An earlier iteration used one policy-controlled agent alongside five
hand-written greedy agents, an architecture common enough to be worth
reporting on. It carried the shared-reward defect quantified in
Section~\ref{sec:credit}, and a second one that is fatal on its own.
Replacing the learned agent with an inert, a uniformly random, a
greedy and the trained agent gives swarm coverage of 82.21\%, 82.37\%,
84.07\% and 84.83\% at 20\% density over 15 seeds. The trained policy
is the best of the four, so it had learned something, but the entire
range spanned by completely different behaviours for that one agent is
1.86 percentage points against a 9.37-point gap to APF. No agent could
have closed it. The hybrid's deficit is a property of its five
reactive agents, not of the learned one, and no conclusion about PPO
follows from it.

The fully decentralised architecture in Section~\ref{sec:ppo}
removes both problems, and the 10.25-point margin between MARL and the
hybrid in Table~\ref{tab:stats} quantifies what the shortcut cost.

\subsection{Repaired Defects}

\textbf{Uncontrolled obstacle seeding.} The Mesa models originally
drew the obstacle seed from Mesa's internal generator, not the seed
supplied by the runner, so one seed gave different layouts on repeat:
503, 510 and 512 obstacle cells across three runs. Both models now
take an explicit seed, verified cell by cell, which is what licenses
the paired tests.

\textbf{Seed Set B was not held out.} The earlier PPO environment
mapped the episode seed with \texttt{s = \_ep\_seed \% 50}, so Set B
resolved to Set A's fifty layouts. The modulo was removed and the sets
are now disjoint. \textbf{Unequal episode budgets.} Boids ran
for 800 steps against 700 for the others; all four now use 700.

\textbf{A retracted ablation figure.} An earlier draft reported the
memory-free policy at 95.0\% after 1.85M steps and concluded the
memory features bought convergence speed alone. The surviving
checkpoint re-evaluates to the 1.00M figure exactly, 89.55\%, and no
checkpoint matching 95.0\% exists. Retrained from 1.00M to 1.85M it
scores $90.0\pm4.2\%$, which is the value in
Table~\ref{tab:ablation}, and the conclusion reverses.

\subsection{Withdrawn and Outstanding}

The sensor-noise experiment is withdrawn. For Boids and APF the
perturbation corrupted the obstacle grid itself, once per episode,
marking roughly 398 of 1{,}991 free cells unreachable and imposing a
mechanical ceiling near 80\% before the controller acted at all. For
the hybrid it entered only the observation of the single learned
agent, which contributes 0.4\% of coverage. The comparison is between
two different interventions and supports no conclusion. A valid
version requires the same corruption, representation and
persistence for every controller.

Redundancy is recorded for the MARL controller but not
retrospectively for the classical ones, so no cross-controller
comparison is made.

%% ===================== DISCUSSION ==============================
\section{DISCUSSION}\label{sec:discussion}

\textbf{Answering the question.} A decentralised multi-agent PPO
policy, in which every UAV learns from a local window, executes on its
own network and exchanges no messages, reaches 90\% coverage at every
obstacle density tested and on layouts it has never seen. APF does so too, on the same layouts and by
the same test; the two learned-free baselines do not. Against a
well-tuned APF baseline it wins clearly in
open terrain, ties at intermediate density, and at the hardest density
trails on the development set but not on the held-out one. The answer to the research question is therefore yes
for the criterion and qualified for the comparison: decentralised
learning matches classical reactive control on this task, and beats it
where the geometry offers least structure.

\textbf{Why the crossover occurs.} The two controllers fail in
opposite regimes, and the crossover in Fig.~\ref{fig:fourway} is the
most informative result here. In an empty grid APF has little to work
with: six agents under mutual repulsion spread evenly and then have
nothing to distinguish one direction from another, so they re-sweep
ground. The learned policy carries a directional signal from its own
belief map and keeps moving towards unexplored space, which is exactly
what an empty arena rewards. As obstacles increase, the situation
narrows. Walls and corridors give APF's repulsion term the structure
it needs, channels agents along streets and decorrelates their
paths, while the learned policy faces a harder navigation problem with
a 7$\times$7 window that cannot see around a building. That the 30\%
result is significant on one seed set and not the other suggests the
remaining gap is small relative to layout-to-layout variation. The classical
controller degrades most gracefully in clutter; the learned one
exploits open space best.

Section~\ref{sec:seedvar} limits how far that reading extends. The
crossover is measured at a fixed budget, and at 1.85M the learned
controller reaches 98.8\% at 20\% density and overtakes APF there on
every training seed. What the main experiment establishes is that the
two are close under equal training, not that clutter imposes a ceiling
the learned controller cannot pass.

\textbf{What the absolute-cell result changes.} No controller improves
in absolute terms as obstacles increase, contrary to what percentage
coverage suggests. This matters beyond the present study: coverage
over free cells is the standard metric across the literature surveyed
in Section~\ref{sec:related}, so any paper varying obstacle density
while reporting only that ratio invites the same misreading.

\textbf{Relation to the literature.} Fusco et al.~\cite{fusco2025}
predict that fully decentralised MARL underperforms under high partial
observability, and that the joint action space makes the design
intractable at scale. This project confirms the second half directly,
having hit the $8^6$ wall, and qualifies the first: with parameter
sharing, a per-agent reward and onboard memory, a fully decentralised
policy is competitive with, and in open terrain better than, a tuned
classical controller. The result suggests the pessimism is about the
naive formulation rather than about decentralisation itself. The
credit-assignment literature~\cite{foerster2018,sunehag2018,rashid2018}
is the reason, and Section~\ref{sec:credit} puts a number on it: the
same network on the shared team reward those papers set out to repair
loses 30.6 points and never reaches the target. The difference between
a policy that fails and one that beats APF was, in this project,
mostly a matter of what the agent was rewarded for.

\textbf{Practical implication.} For open or lightly cluttered search
areas the learned controller covers most ground, and it records
between nine and twenty-six times fewer proximity events than either
classical controller, the margin narrowing as density rises. For dense urban clutter a tuned potential field clears the
same target, remains slightly ahead on coverage and is far cheaper to
deploy, needing no training and no GPU; the learned controller is
worth its cost only if trained past the point where its across-seed
spread closes. And for any mission, adding agents rivals changing algorithm: six to
nine UAVs gains 5.3 to 10.2 points, against an 11.0-point spread
between the best and worst controller at six.

%% ===================== CONCLUSION ==============================
\section{CONCLUSIONS AND FUTURE WORK}\label{sec:conclusion}

A parameter-shared multi-agent PPO policy driving six UAVs from local
observations, with decentralised execution and no message passing,
achieved 90.8\%
to 96.1\% coverage across obstacle densities from 0\% to 30\% on both
development and held-out layouts. Its mean exceeds the 90\%
search-and-rescue planning target in every condition tested, and does
so significantly in seven of eight after correction; APF matches that
record and the other two controllers never reach the target. It beat a tuned APF controller by 4.7 and 5.0
percentage points in open terrain and was statistically
indistinguishable at 10\% and 20\% density; at 30\% it trailed by 3.8
points on the development set, a difference that did not replicate on
the held-out set. It beat a Boids baseline and a single-agent hybrid
by 5.3 to 10.9 points throughout. Transfer to fifty unseen layouts
changed coverage by at most 2.31 points and was never significant.

Two design decisions account for that performance, and the ablations
size both. Replacing the per-agent reward with the shared team reward,
and changing nothing else, costs 30.6 points and leaves none of fifty
runs at the target. Removing the memory features costs a further 4.0 points at one million
steps, and the memory-free policy is still at 90.0\% after 1.85M,
having gained nothing from the extra 850k, so to that budget those
features raise the ceiling and not the speed of reaching it. Neither
decision is a refinement, and they do not fail alike: on the shared
reward the policy collapses to 63.0\% and loses to every baseline
here, while without memory it holds 89.6\%, ahead of Boids and the
hybrid but short of both the target and APF. That second decision also qualifies the headline:
the memory features are information Boids and APF do not have, so the
comparison is between a learned controller with map memory and
classical controllers without it, and on equal information APF wins at
both budgets tested. A third result came from auditing the apparatus:
repeating training at four seeds gives a 7.8-point across-seed spread
at one million steps, which halves by 1.85M where every seed passes
APF and the policy reaches 98.8\% at 20\% density.

\textbf{Contributions.} A four-way comparison on identical matched
layouts with 50 replications, paired significance tests, effect sizes
and bootstrap intervals, none of which is reported by any of the eight
studies surveyed in Table~\ref{tab:related}.
A controlled measurement of the credit-assignment problem, in which
the reward scheme is the only variable and costs more coverage than
the gap between any two controllers here. Reporting of variance across
training runs, which the surveyed literature omits. Evidence that
percentage coverage over free cells misleads when free area varies.
And a working decentralised MARL coverage system with the training
configuration that makes it learnable identified explicitly.

\textbf{Limitations.} The simulation is a 2D kinematic grid with
perfect sensing and localisation, no inertia, wind or battery model,
so only the relative ordering plausibly transfers. The learned
controller carries memory features the classical controllers lack, and
the ablation shows these lift its ceiling, so the headline comparison
is not one of equal information. The main experiment reports one
training seed per condition at a budget the seed study shows to be
undertrained, and that study covers 20\% density only, so the
density-dependent claims are not settled. Swarm size was fixed at six,
hyperparameters were tuned only in the rollout-length dimension, the
sensor-noise experiment is withdrawn, and a separate policy is trained
per density rather than one generalising across them.

\textbf{Future work.} The immediate step is to repeat the four-way
comparison with several training seeds per cell at a budget past
convergence, which the 20\% result suggests would turn parity into an
advantage for the learned controller. Beyond that, one policy trained
across all densities would test generalisation over environment statistics as
well as layouts; minimal communication would test whether the 30\%
deficit is a partial-observability or a policy limit; and a valid
sensor-noise experiment would give the robustness ordering the
withdrawn one could not. Longer term, PyBullet dynamics and Crazyflie
flights at the Bristol Robotics Laboratory would test whether the
ordering survives hardware.

%% ===================== BIBLIOGRAPHY (excluded from count) ======
% References set one step down, as IEEE conference papers normally do.
% All 28 entries are unchanged and complete; only the type size moves.
{\footnotesize\linespread{0.90}\selectfont
\bibliographystyle{IEEEtran}
\bibliography{References}
}

\end{document}
